%% =====================================================================
%%  T-02-05  Persistence Over Attitude
%%  -------------------------------------------------------------------
%%  One idea per file. Core is mandatory. Slots in canonical order.
%%  Max 700 words. No layout commands. No filler. New source material
%%  for this entry goes in sources/<BOOK>/T-02-05.tex -- never by
%%  rewriting this file (docs/RULES.md R0.1).
%% =====================================================================
%% @kind: topic
%% @id: T-02-05
%% @title: Persistence Over Attitude
%% @subtitle: Choose the animal that produces, not the one that looks impressive
%% @chapter: c02
%% @order: 50
%% @status: stable
%% @sources: B1
%% @updated: 2026-09-19
%% @allow-rewrite: slot migration to the seven-question format (docs/RULES.md section 2)

\topic{T-02-05}{Persistence Over Attitude}{Choose the animal that produces, not the one that looks impressive}


\begin{core}
A confident mood is not equipment; it fails exactly when equipment is needed. What survives failure is the decision to keep going, made in advance and executed without reference to how you feel about it.
\end{core}

\begin{mechanism}
Attitude is a state and states are contingent on circumstances, so an attitude that depends on things going well cannot be relied on when they do not. Persistence is a decision, and decisions can be held against feelings. The distinction matters because the two look identical from outside during good periods, which is why people invest in the wrong one. There is also a scale error: morale that works on a group --- a team, a company --- does not transfer to an individual operating alone, and most of what is sold as attitude is group morale repackaged.
\end{mechanism}

\begin{application}
  \item Decide the continuation rule before you start: what you will do after the first, second and third failure.
  \item Judge instruments and people by output, not by presentation.
  \item Treat a bad period as the expected case rather than as evidence that the plan is wrong.
  \item Keep going at reduced scope rather than stopping at full scope.
\end{application}

\begin{feedback}
  \item Your confidence tracks your recent results rather than your plan.
  \item You cannot say what you would do after the third failure, only that you would stay positive.
  \item You are borrowing morale from a group you are not currently part of.
  \item You evaluate options by how impressive they look rather than by what they produce.
\end{feedback}

\begin{failure}
Persistence with no review is just stubbornness, and the source that recommends it also recommends knowing when to walk away, which is the opposite judgment. The rule needs a stopping condition attached to it or it converts sunk cost into virtue. It is also weak equipment for the exchanges in Part IV, where reading and adapting to the counterpart matters more than continuing; persistence is for campaigns, not conversations.
\end{failure}

\begin{countermeasures}
  \item When someone sells you confidence, ask what they do when it runs out. A person with a continuation rule has one; a person with an attitude does not.
  \item Discount impressive presentation in your own assessment of people and options --- the showy animal and the productive animal are priced by different criteria.
  \item Do not let a group's morale stand in for your own plan.
\end{countermeasures}

\sources{B1}
\seesources
\seealso{T-02-03, T-08-02}
